\section{Public Key Cryptography}

We explored the use of private-key cryptography to help Alice and Bob
communicate over an insecure channel.
Assuming they both hold a shared secret key.
They can meet up, or use a trusted third party to deliver the key.
Might work in some cases (e.g. Cold War era “red phone”),
but what if Alice needs to communicate with many more people? Establishing and
storing so many private keys will be problematic.

Enter the Diffie-Hellman Key Exchange Protocol.
In this scheme, they are both capable of 
generating a shared key based on private secret 
keys. It goes as follows:
\begin{itemize}
    \item Alice and Bob agree on group $\mathbb{G}$ of 
    order $m$ and generator $g \in \mathbb{G}$
    (in practice these parameters are standardized).
    \item Alice computes $k_A = P_B^a = (g^b)^a = g^{ab}$
    \item Bob computes $k_B = P_A^bb = (g^a)^b = g^{ab}$
    \item $k_A = k_B$
\end{itemize}
The discrete-logarithm problem guarantees that an 
adversary cannot compute the
secret values of Alice and Bob $(a, b)$
from messages communicated $(P_A, P_B)$.

Diffie and Hellman also introduced the notion of public-key cryptography.
Bob generates a key pair: public key to be disseminated, and secret key to be kept
private. The public key serves as encryption key, 
and secret key serves as decryption key.

Public-key encryption is asymmetric. The receiver generates 
a pair of keys $(pk, sk)$ and the sender uses the public 
key $pk$ to encrypt, the receiever uses the private 
key $sk$ to decrypt. 

The channel must be authenticated or a man in the middle 
may substitute its own public keys and freely decrypt 
in the middle and re-encrypt going the other way. 

Because the adversary has the public key, they may take 
any message and encrypt it. Thus any scheme we create should 
be CPA secure. 